El 6 d'agost de 2012, el robot \textit{Curiosity} va ser dipositat damunt la superfície de Mart per una càpsula d'entrada atmosfèrica ideada pel Mars Science Laboratory. Aquesta càpsula va iniciar l'entrada a l'atmosfera a 125~km de la superfície de Mart i a una velocitat de $5\,845\ \mathrm{m\,s^{-1}}$. Les tècniques usades en el descens van fer que el vehicle arribés a la superfície marciana a una velocitat de només $0,60\ \mathrm{m\,s^{-1}}$. Tenint en compte que la massa del \textit{Curiosity} és de 899~kg, calculeu:

\figura[0.3]{fig1.png}
{\centering\small Imatge del \textit{Curiosity} al cràter Gale, a la superfície de Mart (31 d'octubre de 2012)\par}

\begin{apartat}{1}
L'increment de l'energia mecànica del vehicle en el descens.
\end{apartat}

\begin{apartat}{1}
El mòdul de la intensitat de camp gravitatori que fa Mart en el punt inicial del descens del \textit{Curiosity} i la força (mòdul, direcció i sentit) que el planeta fa sobre el robot en aquest punt.
\end{apartat}

\dades{%
  Massa de Mart, $M_\mathrm{Mart} = 6,42\times 10^{23}\ \mathrm{kg}$.\\
  Radi de Mart, $R_\mathrm{Mart} = 3,39\times 10^{6}\ \mathrm{m}$.\\
  $G = 6,67\times 10^{-11}\ \mathrm{N\,m^2\,kg^{-2}}$.}
